\documentclass[aspectratio=169,10pt]{beamer}

\usetheme{metropolis}
\usepackage[utf8]{inputenc}
\usepackage[T1]{fontenc}
\usepackage{graphicx}
\usepackage{booktabs}
% AI-generated text reused as is goes between slashes (course charter)
\newcommand{\aitext}[1]{/#1/}

\title{Multi-modal hypergraph contrastive learning for medical image segmentation}
\subtitle{Paper presentation -- Modélisation de systèmes intelligents}
\author{Maksym Dolhov \and Hocine Guia}
\institute{M2 Informatique VMI, Université Paris Cité}
\date{02/11/2026}

% Time budget: 15 min total. Speakers: TBD.

\begin{document}

\maketitle

% ~0.5 min | Speaker: TBD
\begin{frame}{Outline}
  % TODO: list the 6 parts of the talk
\end{frame}

% ====================================================================
\section{Motivation}
% ~2 min | Speaker: TBD

\begin{frame}{Context}
  % TODO: multi-modal medical segmentation, why combine modalities
\end{frame}

\begin{frame}{Problem and goal}
  % TODO: few labels, many unlabeled images, self-supervised learning
\end{frame}

% ====================================================================
\section{Related work}
% ~1.5 min | Speaker: TBD

\begin{frame}{Seq SSL vs Joint SSL}
  % TODO: how each works, limits of each (forgetting / conflicts), the gap the paper targets
  % TODO: figure placeholder (cite the paper if reused)
  % \includegraphics[width=\linewidth]{figures/...}
\end{frame}

% ====================================================================
\section{Method}
% ~4.5 min | Speaker: TBD

\begin{frame}{Overview}
  % TODO: the two stages + downstream segmentation, one schema
\end{frame}

\begin{frame}{Stage 1: pretraining}
  % TODO: CycleGAN, cross-modal translation, what it gives (encoder, synthesized images)
\end{frame}

\begin{frame}{Stage 2: fine-tuning}
  % TODO: teacher (hypergraph + contrastive), student (MLP), distillation, pseudo-labels
\end{frame}

\begin{frame}{Segmentation}
  % TODO: U-Net, initialisation, pseudo-labels as supervision
\end{frame}

% ====================================================================
\section{Evaluation}
% ~2.5 min | Speaker: TBD

\begin{frame}{Setup and results}
  % TODO: datasets, % of labels, metrics, comparison, main numbers
\end{frame}

\begin{frame}{Ablation}
  % TODO: which component matters
\end{frame}

% ====================================================================
\section{Conclusion and critique}
% ~2 min | Speaker: TBD

\begin{frame}{Conclusion and critique}
  % TODO: your own judgement (not a summary)
\end{frame}

\begin{frame}{Perspectives}
  % TODO: improvements to consider
\end{frame}

% ====================================================================
\section{Project progress}
% ~2 min | Speaker: TBD

\begin{frame}{Project progress -- Student 1}
  % TODO
\end{frame}

\begin{frame}{Project progress -- Student 2}
  % TODO
\end{frame}

% ====================================================================
\begin{frame}[allowframebreaks]{References}
  \bibliographystyle{plain}
  \nocite{jing2025hycon}
  \bibliography{references}
\end{frame}

\begin{frame}{IAgraphie}
  \IfFileExists{iagraphie.tex}{\input{iagraphie}}{}
\end{frame}

\end{document}
